\begin{textsample}{POS Dim 3 – es – Score 21.00 – es/es\_inf\_10\_camp\_exp.txt}  \label{ex:f3_pos_114}
Campo de Experiências : Traços, \textbf{Sons}, Formas e Cores O Campo de Experiências Traços, \textbf{Sons}, Formas e Cores possibilita à \textbf{criança} desenvolver e valorizar as experiências relacionadas com as diferentes linguagens e manifestações artísticas, culturais, simbólicas e científicas, locais e universais relacionadas aos contextos sociais em que as \textbf{crianças} estão inseridas : família, escola, coletividade.
Tem na observação um modo genuíno de conhecer e interpretar o mundo, valorizando e incentivando a contemplação da vida.
Pensar em proporcionar experiências significa inserir a dimensão da participação das \textbf{crianças} como protagonistas de seu processo de aprendizagem, criando suas próprias produções, isto é, o exercício da autoria ( coletiva e individual ), desenvolvendo, desde muito \textbf{pequenas}, senso estético e crítico diante da realidade que as cerca.
Por meio dessas experiências, as \textbf{crianças} se expressam por várias linguagens e formas de expressão : \textbf{sons}, \textbf{gestos}, formas, traços, encenações, danças, \textbf{mímicas}, modelagens, canções, manipulação de diversos materiais e recursos tecnológicos.
Nesse sentido, na Educação \textbf{Infantil}, é preciso promover o desenvolvimento da sensibilidade, da criatividade, por meio das diversas formas de expressão pessoal e cultural, qualificando e ampliando os repertórios imagéticos, artísticos e simbólicos das \textbf{crianças}, promovendo a abertura ao novo e conhecendo outros modos de expressão e de se relacionar com os objetos de conhecimento, permitindo a elas conhecerem a si mesmas, ao outro e à realidade, potencializando suas singularidades.
No cotidiano da \textbf{instituição} escolar, a organização das diversas experiências deve ocorrer em situações de aprendizagem, sistematicamente estruturadas e propostas pedagógicas bem definidas e intencionais, permitindo às \textbf{crianças} vivenciarem diferentes formas de expressões e linguagens - artes visuais ( \textbf{pintura}, modelagem, fotografia, colagem ), a música, o teatro, a dança e o audiovisual, ampliando seus repertórios culturais.
As DCNEI \textbf{convidam} os educadores a pensar a aprendizagem a partir do que a \textbf{criança} é capaz de fazer, partindo das experiências e de como elas aprendem, levando -as a aprender a partir do que é capaz de fazer. \textbf{Brincar} com tintas, explorar diferentes suportes e materiais, \textbf{brincar}, correr, \textbf{escutar}, \textbf{ouvir} músicas, modelar, recortar, compor, desenhar, fazer uso de instrumentos musicais, criar, encenar, \textbf{dramatizar}, esculpir são situações que possibilitam a expressão e que o professor da Educação \textbf{Infantil} deve proporcionar às \textbf{crianças}, devendo ter clareza de que não pretendemos aqui formar um artista, mas auxiliar, através das diferentes formas de linguagem e da Arte, na construção de seres capazes de expressar sensações, sentimentos, pensamentos e que se tornem potentes para desenvolver seus próprios percursos criativos.
Cabe à escola oferecer contextos de aprendizagem que possam favorecer essas experiências para que todas as \textbf{crianças} desenvolvam seus próprios percursos criativos, os quais são singulares, resultado de sucessivas aprendizagens.
Por isso, as produções das \textbf{crianças} têm valor como parte desses percursos, não requerendo acabamentos, posto que estão em processo.
A produção e sua estética ao olhar do \textbf{adulto} não devem ocupar lugar tão privilegiado quanto os percursos de criação da \textbf{criança}, ou seja, as produções dos \textbf{pequenos} não precisam ser maquiadas ou melhoradas para serem expostas ao público \textbf{adulto}.
É importante levar a \textbf{criança} a perceber a presença da arte no mundo que nos cerca : nas ruas, vitrines, roupas ou na fachada das casas, estimulando o processo criativo e natural das \textbf{crianças}.
É preciso ficar claro para o professor que, quando \textbf{brinca}, a \textbf{criança} desenvolve atividades rítmicas, melódicas, fantasia -se de \textbf{adulto}, produz desenhos, danças, \textbf{inventa} histórias e é esta criatividade natural que deve ser explorada na Educação \textbf{Infantil}.
Na BNCC este Campo de Experiências estabelece que : > O Campo de Experiências “ Traços, \textbf{Sons}, Formas e Cores ” trata dos objetivos de aprendizagem e desenvolvimento, e devem garantir os direitos de aprendizagem de modo a possibilitar à \textbf{criança} : Figura 23 - Direitos de Aprendizagem aplicados ao Campo “ Traços, \textbf{Sons}, Formas e Cores ” ( \textbf{OLIVEIRA}, 2017, p.59 ).
Conviver - e fruir com os \textbf{colegas} e professores manifestações artísticas e culturais da sua comunidade e de outras culturas - artes \textbf{plásticas}, música, dança, teatro, cinema, folguedos e festas populares. \textbf{Brincar} - com diferentes \textbf{sons}, ritmos, formas, cores, \textbf{texturas}, objetos, materiais, construindo cenários e indumentárias para \textbf{brincadeiras} de faz de conta, encenações ou para festas tradicionais.
Participar - de decisões e ações relativas à organização do ambiente ( tanto o cotidiano quanto o preparado para determinados eventos ), à definição de temas e à escolha de materiais a serem usados em atividades \textbf{lúdicas} e artísticas.
Explorar - variadas possibilidades de usos e combinações de materiais, substâncias, objetos e recursos tecnológicos para criar desenhos, modelagens, músicas, danças, encenações teatrais e musicais.
Expressar - suas emoções, sentimentos, necessidades e ideias, cantando, dançando, esculpindo, desenhando, encenando.
Conhecer -se - e construir sua identidade pessoal e cultural, reconhecendo seus interesses na relação com o mundo físico e social.

% matched lemmas: adulto, brincadeira, brincar, colega, convidar, criança, dramatizar, escutar, gesto, infantil, instituição, inventar, lúdico, mímico, oliveira, ouvir, pequeno, pintura, plástico, som, textura
\end{textsample}
